% EXCERPT: original lines 1-64
\section{Spatial moments and reduction to a coordinate direction}
\label{sec:radius}

The finite mean width from Proposition~\ref{prop:mean-width} controls a
regeneration word only in its direction of progress.  We now show that
coexistence of the two escape signs also controls its spatial extent.
This will give two limiting rays and reduce the problem to a coordinate
direction.  The radius argument adapts
\cite[Proposition~3.1]{OpenAI2026}, with a marked connector and a
separated infinite continuation supplying the finite-range version.

For a regeneration word $\gamma$ with position sequence
$(x_0,\ldots,x_m)$, written relative to $x_0=0$, define its radius by
\[
  \mathcal R(\gamma)=\max_{0\le j\le m}|x_j|,
\]
where $|\cdot|$ denotes Euclidean norm.  In particular, the terminal
arrival contributes to the radius.  Distances used for finite-range
independence remain $\ell^\infty$ distances.

\subsection{An integrable spatial radius}

We first record an elementary estimate that controls all initial sums at
once.  Stationarity, rather than independence, is enough for this estimate.
It is the stationary-sequence form of Hopf's maximal ergodic inequality;
see Garsia~\cite{Garsia1965}.  We include a disjoint-interval proof.

\begin{lemma}[A maximal estimate for initial averages]
\label{lem:stationary-maximal}
Let $(Z_i)_{i\ge1}$ be a stationary sequence of nonnegative integrable
random variables.  For every integer $n\ge1$ and every $a>0$,
\[
  \PP\left\{\max_{1\le k\le n}\frac1k\sum_{i=1}^k Z_i>a\right\}
  \le \frac{\EE Z_1}{a}.
\]
\end{lemma}

\begin{proof}
Call an integer $j\ge1$ bad if some interval starting at $j$, of length
at most $n$, has average greater than $a$.  Among the bad starts in
$\{1,\ldots,m\}$, select the leftmost one and a witnessing interval.
Continue by selecting the leftmost bad start strictly to the right of
the selected interval, and repeat.  These disjoint intervals cover all
bad starts in $\{1,\ldots,m\}$ and lie in $\{1,\ldots,m+n\}$.
Their total length is at least the number of those bad starts, while
the sum of the $Z_i$ over them is at least $a$ times their total length.
Consequently
\[
  a\,\#\{j\le m:j\text{ is bad}\}\le\sum_{i=1}^{m+n}Z_i.
\]
Taking expectations, using stationarity, dividing by $m$, and letting
$m\to\infty$ proves the assertion.
\end{proof}

\begin{proposition}[Radius moment under coexistence]
\label{prop:radius-moment}
Let $v\ne0$ satisfy $\|v\|_\infty=1$.  If
$P_0(A_v)>0$ and $P_0(A_{-v})>0$, then the regeneration-word laws of
Proposition~\ref{prop:cut-law} satisfy
\begin{equation}
  E_{\nu_v}\mathcal R+E_{\nu_{-v}}\mathcal R<\infty.
  \label{eq:radius-moment}
\end{equation}
\end{proposition}

